\section{问题描述}
\label{sec:model}
\subsection{系统模型与通信协议}
考虑如下 $n$ 维线性高斯系统，其中 $n\ge1$ 为固定有限维数：
\begin{equation}
 x_{t+1}=Ax_t+w_t,\qquad y_t=x_t+v_t\in\R^n.
 \label{eq:model}
\end{equation}
其中，$x_t$ 为系统状态，$y_t$ 为全状态含噪测量，$w_t$ 和 $v_t$ 分别为过程噪声与测量噪声。
设 $x_0\sim N(m_0,P_0)$，$w_t\sim N(0,Q)$，$v_t\sim N(0,R)$，
$A$ 可逆，$P_0,R,Z\succ0$，$Q\succeq0$。初值、各时刻噪声及触发随机数相互独立。
下文 $I$ 为单位矩阵，$\|\cdot\|_2$ 对向量表示欧氏范数、对矩阵表示诱导谱范数，
$\operatorname{tr}$ 表示矩阵的迹，$\lambda_{\min}$ 表示对称矩阵的最小特征值。

图\ref{fig:protocol}给出了系统的通信与攻击过程。传感器通过可信及时反馈获得接收端实际预测 $m_t^-$。
反馈内容未被篡改，但其中的预测值仍可能含有先前载荷攻击造成的偏差。
传感器以私有随机数 $U_t\sim\operatorname{Unif}[0,1]$ 生成发送事件：
\begin{equation}
 \begin{aligned}
 \gamma_t=0\ &\Longleftrightarrow\ U_t\le g(y_t-m_t^-),\\
 g(z)&=\exp(-z^\top Zz/2).
 \end{aligned}
 \label{eq:trigger}
\end{equation}
$\gamma_t=0$ 表示静默，$\gamma_t=1$ 表示发送。事件位可信，攻击者在发送决定之后修改载荷，接收值为
\begin{equation}
 \widetilde y_t=y_t+a_t,\qquad \|a_t\|_2\le b.
 \label{eq:peak}
\end{equation}
其中，$a_t$ 为载荷修改量，$b$ 为逐路径的欧氏峰值上限，静默时令 $a_t=0$。
本文考虑无丢包、无时延的状态估计过程，估计值不参与物理系统的反馈控制。

\begin{ecustfigure}
\centering
\begin{tikzpicture}[>=Stealth,font=\fontsize{11}{13}\selectfont,
 box/.style={draw,line width=.6pt,rounded corners=2pt,align=center,minimum height=12mm,text width=36mm,inner xsep=2pt},
 edge/.style={->,line width=.7pt},
 edge label/.style={font=\fontsize{10.5}{12.5}\selectfont}]
\node[box] (sensor) {Sensor and\\random trigger};
\node[box,right=18mm of sensor] (attack) {Payload modification\\$\widetilde y_t=y_t+a_t$};
\node[box,right=18mm of attack] (receiver) {Receiver\\$(m_t^+,P_t^+)$};
\node[box,below=20mm of attack] (logger) {Remote logger\\$(P_k^+,\gamma_{k+1})$};
\draw[edge] (sensor) -- node[above,edge label]{sent $y_t$} (attack);
\draw[edge] (attack) -- (receiver);
\draw[edge] (sensor.north) -- ++(0,7mm) -| node[pos=.25,above,edge label]{trusted event $\gamma_t$} (receiver.north);
\draw[edge] (receiver.south) -- ++(0,-39mm) -| node[pos=.25,below,edge label]{trusted actual prediction feedback $m_{t+1}^-$} (sensor.south);
\draw[edge] (sensor.south east) -- ++(3mm,-9mm) -| node[pos=.25,above,edge label]{events} (logger.north west);
\draw[edge] (receiver.south west) -- ++(-3mm,-9mm) coordinate (covturn)
 -- node[midway,above,edge label,xshift=-2mm]{nominal $P_k^+$} (logger.north east |- covturn) -- (logger.north east);
\end{tikzpicture}
\figcaption{实际预测反馈下的载荷攻击与可信记录}{Payload attacks and trusted records with feedback of the receiver's actual prediction}
\label{fig:protocol}
\end{ecustfigure}


\subsection{状态估计与攻击模型}
记 $R_1=R$、$R_0=R+Z^{-1}$。对正定矩阵 $P$，定义
\begin{equation}
 \begin{aligned}
 K_i(P)&=P(P+R_i)^{-1},\\
 F_i(P)&=P-P(P+R_i)^{-1}P\\
       &=(P^{-1}+R_i^{-1})^{-1},\quad i\in\{0,1\}.
 \end{aligned}
 \label{eq:gains}
\end{equation}
接收端按正常模型计算
\begin{equation}
 \begin{aligned}
 m_t^+&=m_t^-+\gamma_tK_1(P_t^-)(\widetilde y_t-m_t^-),\\
 P_t^+&=F_{\gamma_t}(P_t^-),\\
 m_{t+1}^-&=Am_t^+,\quad P_{t+1}^-=AP_t^+A^\top+Q,
 \end{aligned}
 \label{eq:receiver}
\end{equation}
静默时按 $m_t^+=m_t^-$ 更新，发送时使用收到的载荷。
初值为 $(m_0^-,P_0^-)=(m_0,P_0)$。$P_t^\pm$ 是接收端按事件历史递推的名义协方差，
受攻时还需计入估计偏差，才能得到真实均方误差。
对任一有限时刻，全部事件历史对应的 $P_t^\pm$ 均正定，且构成有限值集。
增益矩阵 $K_i$ 一般不对称，后文的范数估计不要求 $\|K_i\|_2\le1$。

用 $S$ 表示独立于物理过程和触发随机数的攻击私有种子，
令 $Y_t^{\rm true}=y_t$ 在发送时成立，静默时为 $\bot$。定义
\begin{equation}
 \begin{aligned}
 \mathcal G_t^-&=\sigma(S,\gamma_{0:t-1},Y_{0:t-1}^{\rm true}),\\
 \mathcal G_t^+&=\mathcal G_t^-\vee\sigma(\gamma_t,Y_t^{\rm true}).
 \end{aligned}
 \label{eq:information}
\end{equation}
攻击策略为任意 $\mathcal G_t^+$ 可测且满足式\eqref{eq:peak}的随机因果策略。
给定策略后，攻击者可由上述历史重建已实施的动作及接收端预测。
攻击者不能读取隐藏状态、未发送测量、未来噪声或触发私有随机数。
式\eqref{eq:information}中的扩展历史用于分析真实状态的条件分布。
远端评估者没有真实已发送测量的副本，仅能利用下文定义的通信记录。

\subsection{攻击损伤与统计隐蔽约束}
固定目标时刻 $k\in\{0,1,\ldots\}$。下标 $0$ 表示无攻击协议的分布，$\pi$ 表示受攻协议的分布。
两种运行均根据各自的实际预测形成发送事件，事件轨迹因而可以不同。
以真实总均方误差（mean-square error，MSE）相对无攻击运行的增量衡量攻击损伤：
\begin{equation}
 J_k(\pi)=\E_\pi\|x_k-m_k^+\|_2^2-
             \E_0\|x_k-m_k^+\|_2^2.
 \label{eq:damage}
\end{equation}
令 $H_t$ 收集截至 $t$ 的可信事件、实际预测反馈及收到的载荷（静默记为 $\bot$）。
日志保存实际发生的事件、反馈和接收值，因此其中的载荷可以包含攻击修改。
为比较不同记录条件下的损伤保证，分别考虑完整日志、事件串及协方差标记记录：
\begin{equation}
 \begin{aligned}
 R_{\rm full}&=H_{k+1},\\
 R_{\rm bits}&=\gamma_{0:k+1},\\
 R_{\rm marked}&=(P_k^+,\gamma_{k+1}).
 \end{aligned}
 \label{eq:records}
\end{equation}
其中，$P_k^+$ 为未经量化的完整协方差矩阵标签。三类记录均包含目标时刻后的下一事件，
标记记录中的协方差用于确定该事件的正常概率基准。

用 $\TV(\mu,\nu)=\sup_E|\mu(E)-\nu(E)|$ 记总变差，并定义
\begin{equation}
 \begin{aligned}
 \Delta_j(\pi)&=\TV(\Law_\pi(R_j),\Law_0(R_j)),\\
 V_j(b,\varepsilon)&=\sup_{\pi:\,\|a_t\|_2\le b,\,\Delta_j(\pi)\le\varepsilon}J_k(\pi).
 \end{aligned}
 \label{eq:value}
\end{equation}
式\eqref{eq:value}以峰值 $b$ 和记录联合分布的TV预算 $\varepsilon$ 同时限制攻击，
$V_j$ 表示满足这两项约束的最坏损伤。该隐蔽约束不依赖具体检测器的报警阈值。
完整日志确定事件串，事件串进一步确定精确标记记录，故数据处理性质给出
$\Delta_{\rm marked}\le\Delta_{\rm bits}\le\Delta_{\rm full}$。

\section{双预算约束下的损伤分析}
\label{sec:bounds}
载荷修改一方面引入估计偏差，另一方面通过实际预测反馈改变后续发送概率。
由于名义协方差由事件历史决定，计算真实损伤时需要同时考虑这两种影响。
为此，先给出条件误差分解，再建立峰值约束与记录约束下的损伤上界，并构造达到相同阶数的攻击。
\begin{lemma}[条件误差分解]
\label{lem:gaussian}
对任意允许策略，$x_t\mid\mathcal G_t^\pm$ 为高斯分布，均值记为 $\mu_t^\pm$，
协方差为 $P_t^\pm$。令 $d_t=\mu_t^--m_t^-$、$D_t=\mu_t^+-m_t^+$，则
\begin{equation}
 \begin{aligned}
 D_t&=(I-K_{\gamma_t}(P_t^-))d_t-\gamma_t K_1(P_t^-)a_t,\\
 d_{t+1}&=AD_t,\quad d_0=0,\\
 J_k(\pi)&=\E_\pi\|D_k\|_2^2+
 \E_\pi\operatorname{tr}P_k^+-\E_0\operatorname{tr}P_k^+.
 \end{aligned}
 \label{eq:decomposition}
\end{equation}
\end{lemma}
\begin{proof}
发送时，给定真实测量后，发送事件的似然因子不再依赖状态，
故在后验归一化时消去，条件分布按标准高斯公式更新。
静默时，将触发概率对测量噪声积分，所得似然正比于
$\exp[-(x_t-m_t^-)^\top R_0^{-1}(x_t-m_t^-)/2]$。
因此静默后的条件均值为 $\mu_t^-+K_0(P_t^-)(m_t^--\mu_t^-)$，
协方差为 $F_0(P_t^-)$。将条件均值与接收端估计相减，得到偏差递推。
$x_t-\mu_t^+$ 在给定 $\mathcal G_t^+$ 后均值为零，
误差平方的交叉项因此具有零期望，得到最后一式。
完整的归纳证明见附录\ref{app:conditional}。
\end{proof}

式\eqref{eq:decomposition}中，第一项为条件偏差引起的损伤，后两项之差来自事件历史概率的变化。
同一事件历史下的协方差递推不变，但不同历史的概率权重可以改变，因此协方差变化项可能为负。
单个策略的 $J_k$ 也可能为负；最坏值 $V_j$ 的可行策略包含零动作策略，故 $V_j\ge0$。
在此基础上，可得以下双预算损伤界。

\Needspace{13\baselineskip}
\begin{theorem}[任意固定维数下的最坏损伤界]
\label{thm:main}
固定式\eqref{eq:model}--\eqref{eq:information}的模型、有限 $k$ 及参考峰值
$b_{\rm ref}>0$。存在 $0<c\le C<\infty$，使所有 $0<b\le b_{\rm ref}$、
$0\le\varepsilon\le1$ 及 $j\in\{\mathrm{full},\mathrm{bits},\mathrm{marked}\}$ 满足
\begin{equation}
 c\min\{b^2,\varepsilon\}\le V_j(b,\varepsilon)
 \le C\min\{b^2,\varepsilon\}.
 \label{eq:main}
\end{equation}
特别地，$V_j(b,0)=0$。常数可依赖维数、模型、$k$ 和 $b_{\rm ref}$，
但不依赖变化的 $b,\varepsilon$。
\end{theorem}

\subsection{峰值约束下的损伤上界}
令 $\mathcal P_t^-$ 为长度 $t$ 的所有事件前缀产生的先验协方差有限值集。
以下最大值均取 $P\in\mathcal P_t^-$、$i\in\{0,1\}$，定义
\begin{equation}
 \begin{aligned}
 \alpha_t&=\max_{P,i}\|I-K_i(P)\|_2,\quad
 \beta_t=\max_P\|K_1(P)\|_2,\\
 r_0&=\beta_0,\quad r_t=\alpha_t\|A\|_2r_{t-1}+\beta_t,\\
 \ell_0&=0,\quad \ell_t=\|A\|_2r_{t-1}\quad(t\ge1).
 \end{aligned}
 \label{eq:radii}
\end{equation}
由引理\ref{lem:gaussian}递推可得 $\|d_t\|_2\le\ell_tb$、$\|D_t\|_2\le r_tb$。
这两个半径界对全部允许策略和事件历史统一成立。
给定事件前缀 $h$，令 $\Sigma_h=P_t^-(h)+R$、
$W_h=(Z^{-1}+\Sigma_h)^{-1}$、
$s_h=\det(I+Z^{1/2}\Sigma_hZ^{1/2})^{-1/2}$。高斯积分给出
\begin{equation}
 \Pr_\pi(\gamma_t=0\mid h)
 =s_h\E_\pi[e^{-d_t^\top W_hd_t/2}\mid h].
 \label{eq:prefix-silence}
\end{equation}
无攻击时的条件静默概率为 $s_h$。由 $W_h\preceq Z$ 及 $1-e^{-u}\le u$，
每步概率差至多为 $\|Z\|_2\ell_t^2b^2/2$。
在共同事件前缀上依次耦合两个伯努利条件分布，得到
\begin{equation}
 \TV(\Law_\pi(\gamma_{0:k}),\Law_0(\gamma_{0:k}))
 \le\frac{\|Z\|_2b^2}{2}\sum_{t=1}^k\ell_t^2.
 \label{eq:prefix-tv}
\end{equation}
记 $\mathcal P_k^+$ 为后验协方差有限值集，
$\Omega_k=\max_{P\in\mathcal P_k^+}\operatorname{tr}P-
\min_{P\in\mathcal P_k^+}\operatorname{tr}P$。
对协方差迹应用有界函数期望差的TV界，再结合式\eqref{eq:decomposition}，得到
\begin{equation}
 J_k\le C_{\rm peak}b^2,\qquad
 C_{\rm peak}=r_k^2+\frac{\Omega_k\|Z\|_2}{2}\sum_{t=1}^k\ell_t^2.
 \label{eq:peak-bound}
\end{equation}

\subsection{记录约束下的损伤上界}
令 $P=P_k^+$、$D=D_k$、$Y=\ind\{\gamma_{k+1}=0\}$，定义
\begin{equation}
 \begin{aligned}
 \Sigma(P)&=APA^\top+Q+R,\\
 s(P)&=\det(I+Z^{1/2}\Sigma(P)Z^{1/2})^{-1/2},\\
 w(P)&=s(P)^{-1},\\
 H(P)&=\tfrac12 A^\top[Z^{-1}+\Sigma(P)]^{-1}A.
 \end{aligned}
 \label{eq:witness-functions}
\end{equation}
目标时刻的偏差 $D$ 经状态转移成为下一时刻的预测偏差 $AD$。
给定 $\mathcal G_k^+$，下一测量相对实际预测服从 $N(AD,\Sigma(P))$，因而
\begin{equation}
 \begin{aligned}
 \E_\pi[Y\mid\mathcal G_k^+]&=s(P)e^{-D^\top H(P)D},\\
 \E_0[w(P)Y]&=1,\quad
 \E_\pi[w(P)Y]=\E_\pi e^{-D^\top H(P)D}.
 \end{aligned}
 \label{eq:witness}
\end{equation}
由式\eqref{eq:witness}可知，当前估计偏差通过反馈降低下一时刻的条件静默概率。
$s(P)$ 为相应协方差下的正常概率，权重 $w(P)$ 将不同事件历史的正常基准统一为1。
因此，加权静默统计量的期望差能够反映当前偏差；下一时刻的载荷修改发生在发送决定之后，不影响这一关系。
$A$ 可逆使全部 $H(P)\succ0$。令
\begin{equation}
 \begin{aligned}
 h_*&=\min_{P\in\mathcal P_k^+}\lambda_{\min}(H(P)),\quad
 M=\max_{P\in\mathcal P_k^+}w(P),\\
 B_*&=r_kb_{\rm ref},\quad q=\frac{1-e^{-h_*B_*^2}}{B_*^2}>0.
 \end{aligned}
 \label{eq:witness-constants}
\end{equation}
凹函数 $1-e^{-h_*x}$ 在 $[0,B_*^2]$ 上的割线给出
$1-e^{-D^\top H(P)D}\ge q\|D\|_2^2$。又 $0\le w(P)Y\le M$，故
\begin{equation}
 \begin{aligned}
 q\E_\pi\|D\|_2^2&\le\E_0[w(P)Y]-\E_\pi[w(P)Y]\\
 &\le M\Delta_{\rm marked}(\pi),\\
 J_k(\pi)&\le C_{\rm record}\Delta_{\rm marked}(\pi),\\
 C_{\rm record}&=M/q+\Omega_k.
 \end{aligned}
 \label{eq:record-bound}
\end{equation}
将式\eqref{eq:peak-bound}与式\eqref{eq:record-bound}结合，并利用三类记录间的数据处理关系，
即可得到定理\ref{thm:main}的上界。该上界同时计入偏差平方与协方差期望变化。

\subsection{可达到损伤下界的攻击构造}
为证明上述损伤阶数可以达到，需要构造满足因果权限与两项预算的攻击。
以下采用私有随机符号，并通过下一拍补偿控制完整日志的联合分布差异。
固定单位方向 $v\in\R^n$，在目标时刻前不施加攻击。
令 $h=\gamma_{0:k-1}$（$k=0$ 时为空串），$K_h=K_1(P_k^-(h))$，并取
\begin{equation}
 \kappa=\bigl[\max\{1,\max_h\|AK_hv\|_2\}\bigr]^{-1},\quad \eta=\kappa b.
 \label{eq:kappa}
\end{equation}
目标时刻发送时注入 $a_k=\sigma\eta v$，其中 $\sigma$ 是私有的等概率正负符号。
若下一时刻也发送，注入 $a_{k+1}=\sigma\eta AK_hv$；
目标时刻静默或其他时刻均不攻击。两个动作都满足欧氏峰值 $b$。
该策略只使用事件历史、私有符号及已知模型，因而满足因果信息约束。

记无攻击运行的前缀概率为 $w_h$、目标时刻的条件发送概率为 $p_h$。
由于目标时刻之前未施加攻击，目标时刻的协方差分布保持不变，故
\begin{equation}
 J_k=L_v\eta^2,\qquad L_v=\sum_hw_hp_h\|K_hv\|_2^2>0.
 \label{eq:lower-damage}
\end{equation}
% Essential joint-log proof; explicit constants are expanded in Appendix B.
由于两拍载荷使用同一私有符号，需要计算其联合分布，才能检验完整日志约束。
给定攻击前完整历史，令
$S_h=P_k^-(h)+R$、$S'_h=AF_1(P_k^-(h))A^\top+Q+R$，
$u_h=AK_hv$，并以 $\phi_\Sigma$ 表示 $N(0,\Sigma)$ 的密度。
当前真实创新的未归一化发送密度为
$f_h(u)=\phi_{S_h}(u)[1-g(u)]$，其积分为 $p_h$。
给定当前真实发送测量，下一真实条件创新 $e$ 服从 $N(0,S'_h)$，
且该条件分布不随当前创新改变。
对有符号幅值 $t$，下一预测偏差为 $-tu_h$；
若下一时刻发送，补偿使接收创新恰为 $e$。
下一记录既可能是静默事件，也可能包含连续载荷。将两者共同表示为
$\{\bot\}\sqcup\R^n$ 上的概率核，并定义当前发送分支的联合核：
\begin{equation}
 \begin{aligned}
 \mathcal K_{h,d}(\{\bot\})&=\int\phi_{S'_h}(e)g(e+d)\,de,\\
 \mathcal K_{h,d}(de)&=\phi_{S'_h}(e)[1-g(e+d)]\,de,\\
 F_h(t)(dz,d\xi)&=f_h(z-tv)\,dz\,\mathcal K_{h,-tu_h}(d\xi).
 \end{aligned}
 \label{eq:main-joint-kernel}
\end{equation}
其中，$z$ 为当前接收创新，$F_h(t)$ 的总质量恒为 $p_h$。
式\eqref{eq:main-joint-kernel}的第一因子描述当前载荷的平移，第二因子描述预测反馈引起的下一事件概率变化。

记 $\|\cdot\|_{\rm var}$ 为有符号测度的变差范数，并令 $a_h=u_h^\top Zu_h$。
由 $\sup_z|\partial_{u_h}g(z)|\le\sqrt{a_h/\mathrm e}$、
$\sup_z|\partial_{u_h}^2g(z)|\le a_h$，
分别计入静默原子和连续密度的导数，得到
\begin{equation}
 \begin{aligned}
 \|F_h''(t)\|_{\rm var}
 &\le \|\partial_v^2 f_h\|_1
 +4\sqrt{a_h/\mathrm e}\|\partial_v f_h\|_1
 +2p_ha_h\\
 &\le B_h<\infty.
 \end{aligned}
 \label{eq:main-kernel-smooth}
\end{equation}
式\eqref{eq:main-kernel-smooth}中的三项依次来自当前密度的二阶导数、
两个因子的交叉导数及下一记录核的二阶导数。
$f_h$ 的高斯方向导数可积，平移保持其 $L^1$ 范数；
在共同支配测度 $dz\otimes(\delta_\bot+d\xi)$ 下，乘积核因而二阶连续 $L^1$ 可微。
固定模型和有限事件前缀使 $B_h$ 与变化的 $b,\varepsilon$ 无关；
附录\ref{app:full-kernel}给出其显式选择。
等概率正负符号使联合核的一阶扰动相消。
对 $Q_h^\eta=[F_h(\eta)+F_h(-\eta)]/2$ 使用中心 Taylor 公式，可得
\begin{equation}
 \begin{aligned}
 &Q_h^\eta-F_h(0)\\
 &\quad=\frac12\int_0^\eta(\eta-t)
 \bigl[F_h''(t)+F_h''(-t)\bigr]\,dt,\\
 &\|Q_h^\eta-F_h(0)\|_{\rm var}\le \tfrac12B_h\eta^2.
 \end{aligned}
 \label{eq:main-kernel-taylor}
\end{equation}
给定攻击前历史及 $(z,\xi)$，按
$m_{k+1}^-=A(m_k^-+K_hz)$，并在 $\xi\ne\bot$ 时令下一载荷为
$m_{k+1}^-+\xi$，即可重建全部事件、载荷和预测记录。
该重建映射不依赖 $t$ 或私有符号；当前静默分支则与无攻击运行相同。
对共同前史积分，利用 TV 在固定映射下的收缩性，以及等质量测度差的
TV 为其变差范数的一半，得到
$\Delta_{\rm full}\le C_\pm\eta^2$，
其中 $C_\pm=\frac14\sum_hw_hB_h\in(0,\infty)$。

再在实验开始前以私有概率
$\lambda=\min\{1,\varepsilon/(C_\pm\kappa^2b^2)\}$ 选择上述攻击，其余情况下不施加攻击，得到
\begin{equation}
 \begin{aligned}
 \Delta_{\rm full}&\le\varepsilon,\\
 J_k&=L_v\min\{\kappa^2b^2,\varepsilon/C_\pm\}.
 \end{aligned}
 \label{eq:attainable}
\end{equation}
取 $c=L_v\min\{\kappa^2,C_\pm^{-1}\}$ 即得全部三类记录的下界，主定理得证。
上述构造中的下一拍动作仅补偿接收创新的平移，不消除条件估计偏差，
也不改变已经形成的目标时刻损伤。它使载荷分布与后续事件概率能够在同一联合核中分析，从而满足完整日志预算。

\begin{remark}[维数、观测条件和时间范围]
式\eqref{eq:main}对任意固定 $n$ 给出同一阶数结构，其中的常数依赖具体维数与模型。
可逆性保证下一时刻事件所对应的二次型在全部偏差方向上正定。
若 $A$ 奇异，则对非零 $v\in\ker A$ 有 $v^\top H(P)v=0$；
沿该方向的当前偏差不改变下一时刻预测均值，因而不能由下一静默概率约束。
这是单步事件感知全部偏差方向所需的结构条件。
若使用部分测量 $y=Cx+v$，对应二次型含有 $CA$，
仅凭 $(A,C)$ 可观不能保证单步二次型正定。
因此，本节结论适用于全状态含噪测量及固定有限时刻；
部分测量情形与时间一致界需要另行分析。
有限时刻分析允许 $A$ 不稳定。
\end{remark}
